% Substitui somente a metodologia. As equacoes agrupam os termos implementados.
\section{Metodologia}
\label{sec,metodologia}

O desenvolvimento foi organizado em etapas de modelagem, integração e experimentação iterativa. Inicialmente, foi preparado o modelo do Atom no Onshape e realizada sua exportação para o formato MJCF. Em paralelo, foram conduzidos testes exploratórios com a infraestrutura de simulação e treinamento disponibilizada por Fabio Gianzanti. Em seguida, o Atom foi integrado a um ambiente de aprendizado por reforço, no qual foram desenvolvidas diferentes formulações de recompensa e treinadas políticas com Proximal Policy Optimization (PPO). O processo evoluiu do incentivo ao deslocamento frontal para a navegação até um alvo e, posteriormente, para o aprendizado com referência periódica de marcha.

\subsection{Modelagem e exportação do Atom}

A estrutura do robô foi organizada no Onshape, com a montagem dos componentes e a definição das articulações. O modelo considerado possui 20 articulações atuadas: seis em cada perna, três em cada braço e duas na cabeça. As propriedades de massa e inércia e as relações cinemáticas entre os corpos foram incorporadas à representação utilizada na simulação.

A exportação foi realizada com a ferramenta \textit{Onshape to Robot}, desenvolvida pela equipe Rhoban \cite{gondry2019,allali2024}, utilizando a saída para MuJoCo. O arquivo \texttt{robot.xml} descreve corpos, articulações, atuadores e geometrias, enquanto \texttt{scene.xml} inclui esse modelo e os elementos do ambiente, como solo e alvo. A configuração de exportação incluiu a combinação e a simplificação de malhas destinadas a colisão. Os parâmetros efetivamente utilizados correspondem aos XMLs carregados em cada experimento, e não necessariamente às opções presentes em arquivos de rascunho.

Na integração foram verificados os nomes e os sentidos das juntas, sua associação com os atuadores e o contato dos pés com o solo. Foram utilizados atuadores de posição, com limites de comando e de força por articulação. Esses parâmetros são escolhas de simulação e não constituem, por si só, uma validação dos limites do hardware real.

\subsection{Infraestrutura de simulação e integração}

Os testes com o projeto de Gianzanti \cite{gianzanti2022} permitiram familiarização com o MuJoCo e com a estrutura de treinamento de locomoção. Essa infraestrutura serviu como ponto de partida para a integração do Atom; portanto, o desenvolvimento não corresponde a uma implementação integralmente independente da base anterior.

A partir dessa base, foram adaptados o carregamento do modelo físico, os nomes das articulações, os comandos dos atuadores, as observações, a recompensa e os critérios de encerramento dos episódios. Também foram desenvolvidos mecanismos para visualização das políticas, registro de experimentos e preservação de modelos acompanhados dos arquivos de ambiente. O ambiente resultante foi registrado no Gymnasium como \texttt{Atom-v1}, com os métodos \texttt{reset()} e \texttt{step()} responsáveis pela reinicialização e pela interação com a simulação.

Os treinamentos de locomoção descritos nesta seção foram executados diretamente pela interface Python do MuJoCo e do Gymnasium. Testes exploratórios em outros ambientes devem ser distinguidos desta etapa: o fluxo atual de treinamento do Atom não depende de comunicação por soquetes com o servidor da RoboCup.

\subsection{Observações, ações e controle}

Na versão com referência de marcha, o vetor de observações possui 62 componentes: 27 posições generalizadas, 26 velocidades generalizadas e nove variáveis adicionais. As posições incluem a translação e o quaternion da base livre, além dos ângulos das 20 juntas. As velocidades incluem as seis componentes da velocidade da base e as velocidades articulares.

As nove variáveis adicionais são o seno e o cosseno da fase da marcha, a altura de cada pé em relação à sua altura inicial, dois indicadores de contato, uma medida de verticalidade do tronco e duas componentes do vetor horizontal até o alvo. As versões anteriores, com outras dimensões de observação, devem ser associadas às respectivas configurações do ambiente.

O espaço de ações possui 20 componentes contínuas em $[-1,1]$. As ações representam correções das referências de posição das juntas, e não torques enviados diretamente pelo agente. Na estratégia com referência, o comando desejado é:
\begin{equation}
\mathbf{u}^{*}_{t}=\operatorname{clip}\!\left(
\mathbf{u}^{\mathrm{ref}}(\phi_t)+\mathbf{S}\mathbf{a}_t\right),
\end{equation}
em que $\mathbf{u}^{\mathrm{ref}}$ é a referência na fase $\phi_t$, $\mathbf{S}$ contém as escalas de correção por atuador e $\mathbf{a}_t$ é a ação da política. O operador de limitação respeita os intervalos de comando. Uma filtragem temporal reduz variações consecutivas:
\begin{equation}
\mathbf{u}_t=\operatorname{clip}\!\left[
\mathbf{u}_{t-1}+\alpha(\mathbf{u}^{*}_t-\mathbf{u}_{t-1})\right].
\end{equation}
As forças resultantes são calculadas pelo simulador conforme os atuadores e a dinâmica do modelo.

\subsection{Treinamento com PPO}

Foi utilizada a implementação de PPO da biblioteca Stable-Baselines3, com política \texttt{MlpPolicy}. No checkpoint protegido adotado como base para os refinamentos, tanto a rede da política quanto a rede de valor possuem duas camadas ocultas com 64 neurônios e ativação Tanh. A política produz a distribuição das ações contínuas, enquanto a rede de valor estima o retorno esperado.

Foram realizados experimentos com diferentes hiperparâmetros e buscas com Optuna. No checkpoint de referência, foram identificados \texttt{n\_steps}=512, \texttt{batch\_size}=64, \texttt{n\_epochs}=10, $\gamma=0{,}98$, $\lambda_{\mathrm{GAE}}=0{,}95$ e coeficiente adicional de entropia igual a zero. Esses valores descrevem essa configuração específica, não todas as tentativas. Nos refinamentos foram empregadas taxas de aprendizado menores, incluindo $10^{-5}$ e $5\times10^{-6}$, buscando limitar a intensidade das atualizações da política previamente treinada.

Os comandos permitiram continuar a partir de checkpoints existentes. Os treinamentos utilizaram PyTorch, com a opção de executar as redes em GPU por CUDA; a simulação física permaneceu sob responsabilidade do MuJoCo. O fluxo atual cria um ambiente de treinamento e outro de avaliação, o que não corresponde a oito ambientes de coleta simultâneos.

\subsection{Evolução da função de recompensa}

\subsubsection*{Deslocamento frontal sem alvo}

As primeiras formulações incentivaram o deslocamento frontal, a manutenção da postura e a continuidade do episódio. Uma versão recuperada do histórico combina o incremento de posição em $x$, uma parcela constante por passo, uma medida de estabilidade e uma penalidade para a magnitude das ações. Em notação agrupada:
\begin{equation}
r_t=w_x\Delta x_t+b+w_sH_t-w_a\|\mathbf{a}_t\|^2
-w_f\mathbb{1}_{\mathrm{queda}}.
\end{equation}
A parcela $H_t$ considera altura e verticalidade do tronco. A expressão resume uma versão inicial específica; sinais e pesos foram modificados ao longo dos experimentos. A penalidade sobre a ação não representa uma medição direta de consumo energético ou de torque.

\subsubsection*{Deslocamento orientado a um alvo}

A tarefa foi reformulada pela inclusão de um alvo fixo na cena. Sua posição horizontal define a direção desejada:
\begin{equation}
\mathbf{e}_t=\frac{\mathbf{p}_{\mathrm{alvo}}-\mathbf{p}_t}
{\|\mathbf{p}_{\mathrm{alvo}}-\mathbf{p}_t\|}.
\end{equation}
Para distância não nula, a velocidade na direção do alvo é a projeção $v^{\mathrm{alvo}}_t=\mathbf{v}_t^{\mathsf T}\mathbf{e}_t$. Também foi empregado o progresso $d_{t-1}-d_t$, com $d_t$ representando a distância horizontal. A distância nula é tratada separadamente para evitar divisão por zero.

As versões orientadas ao alvo incorporaram diferentes combinações de incentivo ao progresso e penalidades para posturas indesejadas, incluindo configurações dos joelhos e orientação corporal. Posteriormente, foram acrescentados termos de alinhamento, desvio lateral, velocidade, projeção do tronco e regularização dos comandos. Não foi utilizada uma única função imutável durante todo o desenvolvimento.

Uma representação geral dessas recompensas é:
\begin{multline}
r_t=b+\sum_i w_iR_i(s_t,a_t)-\sum_j\beta_jP_j(s_t,a_t)\\
+B\mathbb{1}_{\mathrm{chegada}}-F\mathbb{1}_{\mathrm{queda}}.
\end{multline}
Os termos $R_i$ representam incentivos, e $P_j$ representam penalidades. A expressão agrupa componentes por finalidade; as fórmulas e os pesos exatos de cada tentativa são os preservados no código correspondente.

\subsection{Referência de marcha e aprendizado em etapas}

Para orientar a forma do movimento, foi construída uma referência periódica de marcha, denominada \textit{gait}. Essa referência especifica a alternância das pernas e deslocamentos de quadril, joelho e tornozelo, incluindo flexão e intervalos programados de apoio duplo. Uma prévia cinemática permitiu inspecionar as poses e a orientação das solas. Essa prévia altera diretamente as configurações do modelo e não comprova estabilidade dinâmica.

O treinamento seguinte foi dividido em duas etapas. Na primeira, o PPO foi treinado para acompanhar a referência no lugar, com incentivos de reprodução das poses e penalidades para deslocamento horizontal e perda de equilíbrio. Na segunda, a política obtida inicializou o treinamento de deslocamento, incorporando novamente os objetivos de progresso e chegada ao alvo.

A proximidade entre a postura simulada e a referência foi avaliada por um erro normalizado das juntas das pernas:
\begin{gather}
E_t=\frac{1}{|J|}\sum_{j\in J}
\left(\frac{\delta q_{j,t}-\delta q^{\mathrm{ref}}_{j,t}}{\sigma_j}\right)^2,\\
R^{\mathrm{pose}}_t=\exp(-E_t).
\end{gather}
Aqui, $\delta q$ representa o deslocamento em relação à postura inicial e $\sigma_j$ é uma escala por junta. Foram utilizados incentivos à proximidade da referência e penalidades para desvios.

A estratégia constitui um controle residual orientado por referência: parte da estrutura da marcha é definida previamente, e o PPO aprende correções a partir da interação dinâmica e das recompensas. Não se trata de uma caminhada descoberta integralmente do zero, nem de aprendizado por demonstrações humanas. Os contatos efetivos dependem da simulação e podem diferir das fases programadas.

\subsection{Refinamento, avaliação e registro}

Os refinamentos partiram de políticas selecionadas, variando objetivos de velocidade, alinhamento, estabilidade e suavidade. Nas versões recentes, foram acrescentadas medições da velocidade vertical dos pés na aterrissagem e do tempo de contato contínuo do pé de apoio. Elas fundamentaram penalidades suaves para aterrissagens rápidas e para elevação brusca da outra perna antes de uma janela inicial de apoio de 40~ms. Esses termos são incentivos, não restrições que garantam uma sequência de contatos ou ausência de quedas. Sua eficácia deve ser avaliada separadamente.

Na configuração recente, cada ação corresponde a cinco passos físicos de 2~ms, resultando em uma interação a cada 10~ms. O limite de 60~s corresponde a 6.000 interações por episódio. A chegada é reconhecida quando a distância horizontal fica abaixo de 0,30~m e não há queda. O episódio termina por chegada ou queda, ou é truncado no limite temporal. Os critérios de queda consideram altura, verticalidade e inclinação excessiva do tronco; não equivalem apenas a detectar contato de partes do corpo com o solo.

Foram realizadas avaliações periódicas com ações determinísticas, utilizando, em diferentes tentativas, cinco episódios por avaliação. Foram gerados registros de recompensa e duração, curvas de treinamento e avaliação, e checkpoints selecionados pelo maior retorno médio. A inspeção visual no MuJoCo complementou esses registros, permitindo identificar inclinação progressiva, desvio de trajetória e transições bruscas de apoio.

Recompensas de experimentos com formulações diferentes não devem ser interpretadas como uma escala comum de qualidade. Um maior retorno não garante maior naturalidade ou ausência de quedas. A avaliação considera chegada, quedas, tempo de percurso, desvio lateral e características das pisadas. A apresentação de taxas exige contabilização de episódios; uma execução visual isolada não permite afirmar robustez geral.

Foram mantidas cópias identificadas como protegidas, contendo o checkpoint e, quando disponíveis, parâmetros, registros, gráficos, código do ambiente, XMLs e recursos do modelo. Checkpoints intermediários também puderam ser preservados antes da conclusão do treinamento. Quando o snapshot foi capturado posteriormente, essa procedência foi registrada: sua captura não comprova automaticamente identidade com o ambiente existente no início do treino. Assim, distingue-se a preservação dos pesos da preservação completa das condições de execução.
